\documentclass[reqno]{amsart}
\usepackage[T1]{fontenc}
\usepackage{amsmath,amssymb,amsthm}
\usepackage{xr-hyper}
\usepackage[hidelinks]{hyperref}
\externaldocument[ps-]{20261125T010500--poincare-series-and-near-orthogonality__workbook}
\externaldocument[da-]{20261126T030300--the-data__workbook}
\externaldocument[ob-]{20261126T040100--spaces-forms-and-twists__workbook}
\externaldocument[bp-]{20261126T020210--bessel-poincare-series-and-their-gram-operators__workbook}
\externaldocument[m1-]{20261126T051900--unfolding-twice-gram-operators-as-sums-over-rational-points__workbook}
\externaldocument[m3-]{20261126T052200--the-open-cell-siegels-formula-and-poisson-summation__workbook}
\externaldocument[m4-]{20261126T052300--positive-type-and-the-bootstrap-the-second-moment__workbook}

\newcounter{chap}
% Chapter numbers, generated from 20261125T012000--grh-workbook__guide.org
% by scripts/make-chapters.py.  Do not edit: change the order of the
% index and rerun the script (./preview-site.sh and CI do this).
%
% Usage in a sheet, after \newcounter{chap}:
%   % Chapter numbers, generated from 20261125T012000--grh-workbook__guide.org
% by scripts/make-chapters.py.  Do not edit: change the order of the
% index and rerun the script (./preview-site.sh and CI do this).
%
% Usage in a sheet, after \newcounter{chap}:
%   % Chapter numbers, generated from 20261125T012000--grh-workbook__guide.org
% by scripts/make-chapters.py.  Do not edit: change the order of the
% index and rerun the script (./preview-site.sh and CI do this).
%
% Usage in a sheet, after \newcounter{chap}:
%   \input{chapters}\chapnumber{<sheet ID>}
\makeatletter
\providecommand\chapters@set[2]{\@namedef{chapters@#1}{#2}}
\chapters@set{20261125T001000}{1}% zeros-to-ramp-sums
\chapters@set{20261125T001100}{2}% inert-primes-and-estimate-b
\chapters@set{20261125T005000}{3}% galois-fourier-and-averaging
\chapters@set{20261125T005500}{4}% columns-from-a-sphere
\chapters@set{20261125T000000}{5}% slope-method
\chapters@set{20261126T030100}{6}% a-twist-for-each-prime
\chapters@set{20261125T003000}{7}% the-chain-and-the-square-root
\chapters@set{20261125T010000}{8}% rank-lemma-and-zero-estimate
\chapters@set{20261125T010500}{9}% poincare-series-and-near-orthogonality
\chapters@set{20261124T223000}{10}% elliptic-curves-for-the-proof
\chapters@set{20261125T009500}{11}% fock-model-and-theta-intertwiners
\chapters@set{20261124T223100}{12}% complex-multiplication-and-inert-reduction
\chapters@set{20261125T008000}{13}% one-cm-determinant
\chapters@set{20261124T214500}{14}% supersingular-theta-divisibility
\chapters@set{20261124T214600}{15}% the-isotropic-plane
\chapters@set{20261125T013700}{16}% pairings-on-a-formal-group-of-height-two
\chapters@set{20261125T013500}{17}% hecke-nondegeneracy-on-the-tree
\chapters@set{20261126T030200}{18}% the-bound-is-attained
\chapters@set{20261125T009000}{19}% kronecker-limit-and-chowla-selberg
\chapters@set{20261125T002000}{20}% definite-theta-lifts
\chapters@set{20261125T015500}{21}% siegel-weil-in-the-smallest-case
\chapters@set{20261125T002100}{22}% definite-theta-lifts-and-norms
\chapters@set{20261125T016500}{23}% hecke-operators-and-theta-lifts
\chapters@set{20261125T016700}{24}% divisibility-along-the-chain
\chapters@set{20261125T004000}{25}% denominators-and-sha
\chapters@set{20261125T004100}{26}% tate-twists-and-sha
\chapters@set{20261126T020430}{27}% the-minus-part-by-gauss-sums
\chapters@set{20261125T014100}{28}% congruence-modules
\chapters@set{20261125T014500}{29}% eisenstein-congruences-and-selmer-classes
\chapters@set{20261125T014700}{30}% pseudocharacters-and-the-length-bound
\chapters@set{20261125T011000}{31}% galois-representations-black-box
\chapters@set{20261126T030300}{32}% the-data
\chapters@set{20261126T020210}{33}% bessel-poincare-series-and-their-gram-operators
\chapters@set{20261126T030400}{34}% the-seesaw-identity
\chapters@set{20261126T030500}{35}% divisibility-of-the-composite-lift
\chapters@set{20261126T030600}{36}% theorem-a-short-forms
\chapters@set{20261126T030700}{37}% the-determinant-of-a-chain
\chapters@set{20261126T030800}{38}% the-cost-of-one-lift
\chapters@set{20261126T030900}{39}% denominators-of-projections-to-lifts
\chapters@set{20261126T031000}{40}% twists-with-few-congruences
\chapters@set{20261126T031100}{41}% theorem-c-and-the-main-inequality
\providecommand\chapnumber[1]{%
  \@ifundefined{chapters@#1}%
    {\PackageError{chapters}{Sheet #1 is not listed in the index}%
      {List it in 20261125T012000--grh-workbook__guide.org\MessageBreak
       and run scripts/make-chapters.py.}}%
    {\edef\thechap{\@nameuse{chapters@#1}}}}
\makeatother
\chapnumber{<sheet ID>}
\makeatletter
\providecommand\chapters@set[2]{\@namedef{chapters@#1}{#2}}
\chapters@set{20261125T001000}{1}% zeros-to-ramp-sums
\chapters@set{20261125T001100}{2}% inert-primes-and-estimate-b
\chapters@set{20261125T005000}{3}% galois-fourier-and-averaging
\chapters@set{20261125T005500}{4}% columns-from-a-sphere
\chapters@set{20261125T000000}{5}% slope-method
\chapters@set{20261126T030100}{6}% a-twist-for-each-prime
\chapters@set{20261125T003000}{7}% the-chain-and-the-square-root
\chapters@set{20261125T010000}{8}% rank-lemma-and-zero-estimate
\chapters@set{20261125T010500}{9}% poincare-series-and-near-orthogonality
\chapters@set{20261124T223000}{10}% elliptic-curves-for-the-proof
\chapters@set{20261125T009500}{11}% fock-model-and-theta-intertwiners
\chapters@set{20261124T223100}{12}% complex-multiplication-and-inert-reduction
\chapters@set{20261125T008000}{13}% one-cm-determinant
\chapters@set{20261124T214500}{14}% supersingular-theta-divisibility
\chapters@set{20261124T214600}{15}% the-isotropic-plane
\chapters@set{20261125T013700}{16}% pairings-on-a-formal-group-of-height-two
\chapters@set{20261125T013500}{17}% hecke-nondegeneracy-on-the-tree
\chapters@set{20261126T030200}{18}% the-bound-is-attained
\chapters@set{20261125T009000}{19}% kronecker-limit-and-chowla-selberg
\chapters@set{20261125T002000}{20}% definite-theta-lifts
\chapters@set{20261125T015500}{21}% siegel-weil-in-the-smallest-case
\chapters@set{20261125T002100}{22}% definite-theta-lifts-and-norms
\chapters@set{20261125T016500}{23}% hecke-operators-and-theta-lifts
\chapters@set{20261125T016700}{24}% divisibility-along-the-chain
\chapters@set{20261125T004000}{25}% denominators-and-sha
\chapters@set{20261125T004100}{26}% tate-twists-and-sha
\chapters@set{20261126T020430}{27}% the-minus-part-by-gauss-sums
\chapters@set{20261125T014100}{28}% congruence-modules
\chapters@set{20261125T014500}{29}% eisenstein-congruences-and-selmer-classes
\chapters@set{20261125T014700}{30}% pseudocharacters-and-the-length-bound
\chapters@set{20261125T011000}{31}% galois-representations-black-box
\chapters@set{20261126T030300}{32}% the-data
\chapters@set{20261126T020210}{33}% bessel-poincare-series-and-their-gram-operators
\chapters@set{20261126T030400}{34}% the-seesaw-identity
\chapters@set{20261126T030500}{35}% divisibility-of-the-composite-lift
\chapters@set{20261126T030600}{36}% theorem-a-short-forms
\chapters@set{20261126T030700}{37}% the-determinant-of-a-chain
\chapters@set{20261126T030800}{38}% the-cost-of-one-lift
\chapters@set{20261126T030900}{39}% denominators-of-projections-to-lifts
\chapters@set{20261126T031000}{40}% twists-with-few-congruences
\chapters@set{20261126T031100}{41}% theorem-c-and-the-main-inequality
\providecommand\chapnumber[1]{%
  \@ifundefined{chapters@#1}%
    {\PackageError{chapters}{Sheet #1 is not listed in the index}%
      {List it in 20261125T012000--grh-workbook__guide.org\MessageBreak
       and run scripts/make-chapters.py.}}%
    {\edef\thechap{\@nameuse{chapters@#1}}}}
\makeatother
\chapnumber{<sheet ID>}
\makeatletter
\providecommand\chapters@set[2]{\@namedef{chapters@#1}{#2}}
\chapters@set{20261125T001000}{1}% zeros-to-ramp-sums
\chapters@set{20261125T001100}{2}% inert-primes-and-estimate-b
\chapters@set{20261125T005000}{3}% galois-fourier-and-averaging
\chapters@set{20261125T005500}{4}% columns-from-a-sphere
\chapters@set{20261125T000000}{5}% slope-method
\chapters@set{20261126T030100}{6}% a-twist-for-each-prime
\chapters@set{20261125T003000}{7}% the-chain-and-the-square-root
\chapters@set{20261125T010000}{8}% rank-lemma-and-zero-estimate
\chapters@set{20261125T010500}{9}% poincare-series-and-near-orthogonality
\chapters@set{20261124T223000}{10}% elliptic-curves-for-the-proof
\chapters@set{20261125T009500}{11}% fock-model-and-theta-intertwiners
\chapters@set{20261124T223100}{12}% complex-multiplication-and-inert-reduction
\chapters@set{20261125T008000}{13}% one-cm-determinant
\chapters@set{20261124T214500}{14}% supersingular-theta-divisibility
\chapters@set{20261124T214600}{15}% the-isotropic-plane
\chapters@set{20261125T013700}{16}% pairings-on-a-formal-group-of-height-two
\chapters@set{20261125T013500}{17}% hecke-nondegeneracy-on-the-tree
\chapters@set{20261126T030200}{18}% the-bound-is-attained
\chapters@set{20261125T009000}{19}% kronecker-limit-and-chowla-selberg
\chapters@set{20261125T002000}{20}% definite-theta-lifts
\chapters@set{20261125T015500}{21}% siegel-weil-in-the-smallest-case
\chapters@set{20261125T002100}{22}% definite-theta-lifts-and-norms
\chapters@set{20261125T016500}{23}% hecke-operators-and-theta-lifts
\chapters@set{20261125T016700}{24}% divisibility-along-the-chain
\chapters@set{20261125T004000}{25}% denominators-and-sha
\chapters@set{20261125T004100}{26}% tate-twists-and-sha
\chapters@set{20261126T020430}{27}% the-minus-part-by-gauss-sums
\chapters@set{20261125T014100}{28}% congruence-modules
\chapters@set{20261125T014500}{29}% eisenstein-congruences-and-selmer-classes
\chapters@set{20261125T014700}{30}% pseudocharacters-and-the-length-bound
\chapters@set{20261125T011000}{31}% galois-representations-black-box
\chapters@set{20261126T030300}{32}% the-data
\chapters@set{20261126T020210}{33}% bessel-poincare-series-and-their-gram-operators
\chapters@set{20261126T030400}{34}% the-seesaw-identity
\chapters@set{20261126T030500}{35}% divisibility-of-the-composite-lift
\chapters@set{20261126T030600}{36}% theorem-a-short-forms
\chapters@set{20261126T030700}{37}% the-determinant-of-a-chain
\chapters@set{20261126T030800}{38}% the-cost-of-one-lift
\chapters@set{20261126T030900}{39}% denominators-of-projections-to-lifts
\chapters@set{20261126T031000}{40}% twists-with-few-congruences
\chapters@set{20261126T031100}{41}% theorem-c-and-the-main-inequality
\providecommand\chapnumber[1]{%
  \@ifundefined{chapters@#1}%
    {\PackageError{chapters}{Sheet #1 is not listed in the index}%
      {List it in 20261125T012000--grh-workbook__guide.org\MessageBreak
       and run scripts/make-chapters.py.}}%
    {\edef\thechap{\@nameuse{chapters@#1}}}}
\makeatother

\ifcsname chapters@20261126T052000\endcsname
  \chapnumber{20261126T052000}
\else
  \renewcommand{\thechap}{?}
\fi
\renewcommand{\thesection}{\thechap.\arabic{section}}
\theoremstyle{definition}
\newtheorem{problem}{Problem}[section]
\newtheorem{definition}{Definition}[section]
\theoremstyle{plain}
\newtheorem{theorem}{Theorem}[section]
\theoremstyle{remark}
\newtheorem*{leanremark}{In the Lean formalization}
\newcommand{\Z}{{\mathbb{Z}}}
\newcommand{\C}{{\mathbb{C}}}
\newcommand{\F}{{\mathbb{F}}}
\newcommand{\A}{\mathbb A}
\newcommand{\Oc}{\mathcal O}
\newcommand{\Uc}{{\mathcal{U}}}
\newcommand{\Kc}{\mathcal K}
\newcommand{\Pc}{\mathcal P}
\newcommand{\Bc}{\mathcal B}
\newcommand{\cMc}{\mathcal M}
\newcommand{\Gb}{\bar G}
\newcommand{\Rb}{\bar R_0}
\newcommand{\be}{\mathbf e}
\newcommand{\Sf}{\mathfrak S}
\newcommand{\rB}{\rho_{\mathrm{Bes}}}
\newcommand{\Amp}{\mathrm{Amp}}
\newcommand{\Did}{D_{\mathrm{id}}}
\newcommand{\Dmain}{D_{\mathrm{main}}}
\newcommand{\disc}{\operatorname{disc}}
\newcommand{\lk}{\operatorname{lk}}
\newcommand{\mt}{\operatorname{mt}}
\newcommand{\Fl}{\mathrm{Fl}}
\newcommand{\GL}{\mathrm{GL}}
\newcommand{\diag}{\operatorname{diag}}
\newcommand{\tl}[1]{\texttt{#1}}
\setlength{\emergencystretch}{3em}

\title{Which terms survive: characters, minors and the gain at the level}
\author{}
\date{November 26, 2026}

\begin{document}
\maketitle
\refstepcounter{chap}\label{chap}

\paragraph{Prerequisites.}
From Chapter~\ref{m1-chap}: amplitudes and amplitude sums \(\Amp(\Sf;c,s,c')\) over sets \(\Sf\) of rational points (Definition~\ref{m1-def:amplitude}), the two sizes \(\Did\), \(\Dmain\) (Definition~\ref{m1-def:two-sizes}), the block identity (Theorem~\ref{m1-thm:blocks}), the full Bessel function \(W\) and its cut \(W^\flat\) (Theorem~\ref{m1-thm:bessel-system}), and Problems~\ref{m1-pr:near-one} and~\ref{m1-pr:P-one}.  From Chapter~\ref{bp-chap}: the decomposition \(Y_0=\bigoplus_aEv_a\oplus X_0\oplus Ez\oplus\bigoplus_aEv'_a\) with \(h(z,z)=-N\), the groups \(\Uc\) and \(R=\Uc\rtimes U(X_0)\) (Section~\ref{bp-sec:subgroup}); the coordinate space with basis \(e_{c,i}\), the scaled Poincar\'e map \(\Pc^\circ\), the Gram operators \(\Kc_s\) and the covolume \(\nu\) (Definition~\ref{bp-def:scaled}); \(\Rb\) and \(\rB\) (Definition~\ref{bp-def:bessel-rep}); the lattice basis, in which integrality and reduction are taken.  From Chapter~\ref{ob-chap}: \(\cMc_i\), \(\Gb=\Gb_0\times\Gb_1\) and Problem~\ref{ob-pr:group-order}; from Definition~\ref{da-def:data}: \(b^{10}\le p_i\le b^{11}\), \(M=p_0p_1\); and \(b\ge1000\) (Chapter~\ref{bp-chap}).  From Chapter~\ref{ps-chap}: Theorem~\ref{ps-thm:petersson} and Problem~\ref{ps-pr:J-bound}.  The Bruhat decomposition of \(\GL_n(\F_p)\) \cite[\S8.2]{Ca72}.  Hermitian forms and inner products are linear in the first variable; norms of operators are Hilbert--Schmidt norms.

\medskip
This sheet serves the second moment in the proof of Theorem~\ref{bp-thm:B} of Chapter~\ref{bp-chap}, \emph{Theorem~B: enough forms}.  Using the block identity and the size comparison for the data, this sheet derives the strip estimate of Theorem~\ref{thm:strip} and the estimates at the identity of Problem~\ref{pr:near-one-data} from three results of the proof.  By Theorem~\ref{m1-thm:blocks}, the \(((c,i),(c',i'))\) entry of \(\Kc_s\) is the complex conjugate of \(\nu^{-1}\Amp(U(Y_0)(F);c,s,c')_{ii'}\), a sum over the rational points \(\gamma\) for which \(h_c^{-1}\gamma_fh_{c'}\) is integral with reduction \(s\).  Which sets of terms have amplitude sum zero, and how large is the second moment of the terms outside \(R(F)\) when \(s\) agrees with an element of \(\Rb\) at one level prime?

\section{Characters that disagree}\label{sec:characters}

\begin{problem}\label{pr:agree}
Let \(U,U'\) be subgroups of a finite group \(H\) with characters \(\psi,\psi'\), and let \(\phi\colon H\to\C\) satisfy \(\phi(ugu')=\psi(u)\phi(g)\psi'(u')\) for \(u\in U\), \(u'\in U'\).  If \(ugu'=g\) for some \(u,u'\) with \(\psi(u)\psi'(u')\ne1\), then \(\phi(g)=0\).
\end{problem}

Let \(G=\GL_n(\F_p)\), and let \(U\subset G\) be the upper unitriangular group, \(\psi\) a nontrivial character of \(\F_p\), and \(\psi_U(u):=\psi(\sum_i u_{i,i+1})\).  A double coset \(UgU\) \emph{supports} if \(\psi_U(u)\psi_U(u')=1\) whenever \(ugu'=g\) (this does not depend on the representative \(g\)); by Problem~\ref{pr:agree} the functions \(\phi\) with \(\phi(ugu')=\psi_U(u)\phi(g)\psi_U(u')\) vanish off the supporting double cosets.  By the Bruhat decomposition every double coset is \(UwtU\) with \(w\) a permutation matrix and \(t\) diagonal.

\begin{problem}\label{pr:gl2-cells}
In \(\GL_2(\F_p)\), \(U\diag(t_1,t_2)U\) supports if and only if \(t_1=t_2\), and every double coset in the big cell supports.  So \((p-1)+(p-1)^2=p(p-1)\) of the \(2(p-1)^2\) double cosets support: \(6\) of \(8\) for \(p=3\), \(20\) of \(32\) for \(p=5\).
\end{problem}

\begin{problem}\label{pr:gl3-cells}
By machine: in \(\GL_3(\F_2)\), \(4\) of the \(6\) double cosets support, of sizes \(8,32,32,64\) (the two simple reflections do not); in \(\GL_3(\F_3)\), \(18=2+8+8\) of the \(48\) support, of sizes \(27\), \(243\), \(729\) respectively.
\end{problem}

In all four cases the number of supporting double cosets is \(p^{n-1}(p-1)\).  By Mackey's formula it is \(\dim\operatorname{End}_G(\operatorname{Ind}_U^G\psi_U)\), and since the Gelfand--Graev representation \(\operatorname{Ind}_U^G\psi_U\) is multiplicity free \cite[\S8.1]{Ca85}, it is the number of its irreducible constituents.

\begin{problem}\label{pr:gl2-frame}
On \(\F_p^2\) put \(h(y,y'):=y_1y'_2+y_2y'_1\), and for \(g\in\GL_2(\F_p)\) let \(\langle g\rangle_1:=h(g\be_1,\be_1)\).  Then \(\langle g\rangle_1=g_{21}\); \(\F_pg\be_1=\F_p\be_1\) if and only if \(\langle g\rangle_1=0\); and the \(g\) with \(\langle g\rangle_1=0\) whose double coset supports are exactly the scalar multiples of elements of \(U\).
\end{problem}

\section{Frames, discriminants and links}\label{sec:frames}

In the proof the roles of \(U\), \(\psi_U\) and \(\be_1\) are played by \(\Uc\), the character \(\psi_{\Uc}\) and the flag of \(v_1,\dots,v_\varrho\).  The group \(U(Y_0)(F)\) of rational points contains the rational Bessel subgroup \(R(F)\).

\begin{definition}\label{def:frame}
The \emph{frame} is \((v_1,\dots,v_\varrho,v_{\varrho+1})\) with \(v_{\varrho+1}:=z\).  For \(\gamma\in U(Y_0)(F)\) the matrix of \emph{frame products} is \(\bigl(h(\gamma v_b,v_a)\bigr)_{1\le a,b\le\varrho+1}\): for \(a\le\varrho\) its entry is the coefficient of \(v'_a\) in \(\gamma v_b\), and for \(a=\varrho+1\) it is \(-N\) times the coefficient of \(z\).  Its leading minors are \(\langle\gamma\rangle_j:=\det\bigl(h(\gamma v_b,v_a)\bigr)_{1\le a,b\le j}\) for \(1\le j\le\varrho+1\), and \(\langle\gamma\rangle_0:=1\).
\end{definition}

\begin{definition}\label{def:disc}
The \emph{discriminant} of \(\gamma\) is \(\disc(\gamma):=\langle\gamma\rangle_{\varrho+1}\overline{\langle\gamma\rangle_{\varrho+1}}-N^2\langle\gamma\rangle_\varrho\overline{\langle\gamma\rangle_\varrho}\in E\).
\end{definition}

In this sheet \(\Fl:=Ev_1+\dots+Ev_\varrho\), and for \(y\in Y_0\) we write \([y]\in X_0\oplus Ez\) for its component in \(X_0\oplus Ez\); for \(y\perp\Fl\) it represents the class of \(y\) in \(\Fl^\perp/\Fl\).

\begin{definition}\label{def:link}
The \emph{link} of \(\gamma\) is the subspace \(\lk(\gamma):=\{([y],[\gamma^{-1}y]):y\perp\Fl,\ y\perp\gamma\Fl\}\) of \((X_0\oplus Ez)^2\).
\end{definition}

The \emph{characters agree} on \(\gamma\) if, for all \(u,u'\in\Uc(F)\) with \(u\gamma=\gamma u'\), the character that the proof fixes on the finite adelic points of \(\Uc\) takes the same value at the finite components of \(u\) and \(u'\).  In this sheet
\[
 \Gamma^*:=\{\gamma:\disc(\gamma)\ne0\}\cup\{\gamma:\text{the characters agree on }\gamma,\ \disc(\gamma)=0,\ \gamma\Fl\ne\Fl\}
\]
is the set of \emph{remaining terms}.

\begin{definition}\label{def:scalar-coset}
For \(\zeta\in\Oc_E\) with \(\zeta\bar\zeta=1\), the \emph{scalar coset} \(\zeta\Rb\) is the set of \(s\in\Gb\) for which some \(t\in\Rb\) has \(s_w=(\zeta\bmod w)\,t_w\) for every \(w\in\cMc\).  A set \(\Bc\subseteq\Gb\) \emph{avoids the scalar cosets} if it meets no \(\zeta\Rb\) with \(\zeta\ne1\).
\end{definition}

\begin{theorem}[the terms that survive]\label{thm:survivors}
If \(\Bc\subseteq\Gb\) avoids the scalar cosets and \(s\in\Bc\), then for all classes \(c,c'\), for \(W\) and for \(W^\flat\),
\[
 \Amp\bigl(U(Y_0)(F)\smallsetminus R(F);c,s,c'\bigr)=\Amp(\Gamma^*;c,s,c').
\]
\end{theorem}

\noindent\emph{Status}: read from the proof, where it is stated for an admissible filling (Definition~\ref{bp-def:filling}) over a CM field equipped with a CM type and containing a primitive \(5040\)-th root of unity, which the data provide.  Problem~\ref{pr:agree} is a finite model of one step only, the cancellation of the terms whose characters disagree.  The declaration for the theorem is named in the remark at the end of the chapter.

\begin{problem}\label{pr:two-links}
In a three-dimensional model over \(E\), let \(\varrho=1\) and \(X_0=0\), with basis \(v_1,z,v'_1\), \(h(v_1,v'_1)=1\), \(h(v_1,v_1)=h(v'_1,v'_1)=0\), \(h(z,z)=-N\) and \(z\perp v_1,v'_1\).  The two maps \(\gamma_\pm\) exchanging \(v_1\) and \(v'_1\) with \(z\mapsto\pm z\) are isometries of square \(1\) with \(\langle\gamma_\pm\rangle_1=1\), \(\langle\gamma_\pm\rangle_2=\mp N\), \(\disc(\gamma_\pm)=0\) and \(\gamma_\pm\Fl\ne\Fl\); the characters agree on them, whatever the character of \(\Uc\), so both lie in \(\Gamma^*\), and \(\lk(\gamma_\pm)\) is spanned by \((z,\pm z)\).
\end{problem}

Reductions act on the lattice modulo \(w\) in the lattice basis, in which the frame vectors are basis vectors.  For \(s\in\Gb\) and \(w\in\cMc\), let \(s\) act modulo \(w\) by \(s_w\), and modulo \(\bar w\) by the adjoint inverse of \(s_w\) for the bilinear form whose matrix is the Gram matrix of the lattice basis reduced modulo \(w\), which is invertible.  For \(w\in\cMc_i\), \(\#k_w=p_i\) (Problem~\ref{ob-pr:group-order}).

\begin{definition}\label{def:meet}
The element \(s\) \emph{meets} at \(w\) (or at \(\bar w\)) if, modulo that place, the span of the frame and the span of its image under \(s\) share a nonzero vector.  The \emph{meet norm} \(\mt(s)\) is the product, over \(w\in\cMc\), of a factor \(\#k_w\) if \(s\) meets at \(w\) and a further factor \(\#k_w\) if \(s\) meets at \(\bar w\).
\end{definition}

\begin{definition}\label{def:open}
The element \(s\) is \emph{open} if at some \(w\in\cMc_0\) the leading minors of orders \(1,\dots,\varrho\) of the matrix \(\bigl(h(s_wv_b,v_a)\bigr)_{a,b\le\varrho+1}\), formed with the reduced form, are all nonzero.
\end{definition}

\begin{definition}\label{def:fixer}
Let \(\mathcal D\subseteq\cMc_0\), and for each \(w\in\mathcal D\) choose \(w\) or \(\bar w\) and a nonzero vector \(y_w\) of the span of the frame modulo it.  The \emph{fixer} of these vectors is the subgroup of the \(s\in\Gb\) whose action modulo the chosen place fixes \(y_w\), for every \(w\in\mathcal D\).
\end{definition}

For \(j=0,1\), the \emph{strip} \(\Bc_j\) is the set of \(s\in\Gb\) whose component in \(\Gb_j\) equals that of an element of \(\Rb\).

\begin{problem}\label{pr:scalar-one}
If \(\zeta\in\Oc_E\), \(\zeta\bar\zeta=1\) and \(\zeta\equiv1\) modulo every \(w\in\cMc_j\), then \(\zeta=1\).  For example, in \(\Z[\mathrm i]\) reduction modulo \(2+\mathrm i\) is injective on \(\mu_4\).
\end{problem}

\begin{problem}\label{pr:strip-meets}
Every \(s\in\Bc_j\) meets at \(w\) and at \(\bar w\) for every \(w\in\cMc_j\), so \(\mt(s)\ge p_j^{2f}\ge(M/b)^f\); and \(\Bc_j\) avoids the scalar cosets.
\end{problem}

\begin{problem}\label{pr:fixer-meets}
Every fixer contains the reduction of the finite adelic points of \(U(X_0)\) that preserve the lattice, and each element of a fixer meets at the chosen place over every \(w\in\mathcal D\).
\end{problem}

\section{The gain at the level}\label{sec:gain}

On the strip \(\Bc_j\) every reduction meets at the \(2f\) places \(w,\bar w\), \(w\in\cMc_j\).  In Petersson's formula the same mechanism appears when the terms are sorted by their reduction modulo \(p\).

\begin{problem}\label{pr:petersson-gain}
Let \(p\) be a prime and \(k'\ge2\), \(m',n'\ge1\) integers (primed, since \(k\) and \(m\) belong to the data); for even \(k'\ge4\) these are the weight and the indices of Theorem~\ref{ps-thm:petersson}.  If \(\gamma\in\mathrm{SL}_2(\Z)\) is upper triangular modulo \(p\) and \(\gamma_{21}\ne0\), then \(\lvert\gamma_{21}\rvert\ge p\) and \(\lvert J_{k'-1}(4\pi\sqrt{m'n'}/\lvert\gamma_{21}\rvert)\rvert\le(2\pi\sqrt{m'n'}/p)^{k'-1}/(k'-1)!\).
\end{problem}

\begin{problem}\label{pr:columns}
Let \(p\) be a prime and let \(\be_1,\be_2\) be the standard basis of \(\Z^2\).  The integral columns \(\be_1\) and \(\be_1+p\be_2\) agree modulo \(p\); the lattice they span has determinant \(p\) and Gram determinant \(p^2\) (\(2,4\) for \(p=2\); \(3,9\) for \(p=3\)).
\end{problem}

In this sheet \(\Amp^\flat\) denotes the amplitude sums formed with \(W^\flat\) in place of \(W\), and \(\lVert\cdot\rVert\) the Euclidean norm of the entries.

\begin{definition}\label{def:moments}
For \(\Sf\subseteq U(Y_0)(F)\) and \(\Bc\subseteq\Gb\), the \emph{second moment} of \(\Sf\) over \(\Bc\) is \(\nu^{-2}\sum_{c,c'}\sum_{s\in\Bc}\lVert\Amp(\Sf;c,s,c')\rVert^2\) (and likewise with \(W^\flat\); without qualification, the second moment is formed with \(W\)).  The \emph{link second moment} of \(\Sf\) over \(\Bc\) is \(\sum_{c,c'}\sum_{s\in\Bc}\sum_\Lambda\lVert\Amp^\flat(\Sf\cap\{\lk=\Lambda\};c,s,c')\rVert^2\in[0,\infty]\), the inner sum over the subspaces \(\Lambda\) of \((X_0\oplus Ez)^2\).
\end{definition}

By Theorem~\ref{m1-thm:blocks}, for all classes \(c,c'\), Bessel coordinates \(i,i'\) and \(s\in\Gb\),
\[
 \langle\Pc^\circ e_{c,i},\rho(s)\Pc^\circ e_{c',i'}\rangle=\nu^{-1}\Amp\bigl(U(Y_0)(F);c,s,c'\bigr)_{ii'}.
\]
By Definition~\ref{bp-def:bessel-rep}, \(\langle e_{c,i},\rB(s)e_{c',i'}\rangle=\nu^{-1}\Amp(R(F);c,s,c')_{ii'}\).

\begin{problem}\label{pr:hs-blocks}
\(\lVert\Kc_s-\rB(s)\rVert^2=\nu^{-2}\sum_{c,c'}\lVert\Amp(U(Y_0)(F)\smallsetminus R(F);c,s,c')\rVert^2\), the second moment of \(U(Y_0)(F)\smallsetminus R(F)\) over \(\{s\}\).
\end{problem}

\begin{problem}\label{pr:cross}
For vectors \(q_\Lambda\) of an inner product space, finitely many nonzero, \(\lVert\sum_\Lambda q_\Lambda\rVert^2\le\sum_\Lambda\lVert q_\Lambda\rVert^2+\sum_{\Lambda\ne\Lambda'}\lvert\langle q_\Lambda,q_{\Lambda'}\rangle\rvert\).
\end{problem}

\begin{theorem}[links at meeting reductions]\label{thm:link-bound}
If \(z_0\ge1\) is an integer with \(z_0\le\mt(s)\) for every \(s\in\Bc\), then the link second moment of \(\{\gamma:\disc(\gamma)\ne0\}\) over \(\Bc\) is at most \((b^{2f}/z_0)^k\Dmain\,\nu^2\).
\end{theorem}

\begin{theorem}[the cut]\label{thm:cut}
If \(\Bc\) avoids the scalar cosets, \(C_0\ge0\), and the link second moment of \(\{\gamma:\disc(\gamma)\ne0\}\) over \(\Bc\) is at most \(C_0\nu^2\), then the second moment of \(U(Y_0)(F)\smallsetminus R(F)\) over \(\Bc\) has square root at most \(\sqrt{C_0}+b^{-15}\sqrt{\Did}\).
\end{theorem}

\noindent\emph{Status}: both are read from the proof, where they are stated for an admissible filling over a CM field equipped with a CM type and containing a primitive \(5040\)-th root of unity.  By Theorem~\ref{thm:survivors}, the blocks in Theorem~\ref{thm:cut} are those of \(\Gamma^*\).  In the proof of Theorem~\ref{thm:cut}, the terms of zero discriminant in a cut block form one piece, and the terms of nonzero discriminant form one piece for each link; the squares of the pieces of nonzero discriminant, summed over the blocks, give the link second moment.  The proof of Theorem~\ref{thm:link-bound} uses Hua's determinant integral over complex matrices \cite[Ch.~II]{Hua}, sums over sublattices and Liouville's inequality for Gram determinants; the next two problems give the one-row integral and a sublattice sum over \(\Z^2\).  For related lattice counts see \cite{Sch}.

\begin{problem}\label{pr:hua}
For every integer \(k\ge1\), \(\int_{\C^k}(1+\lvert y\rvert^2)^{-k-1}\,dy=\pi^k/k!\), for Lebesgue measure on \(\C^k=\mathbb R^{2k}\); the values for \(k=1,2,3,6,7\) are approximately \(3.14\), \(4.93\), \(5.17\), \(1.34\), \(0.60\), and the integral is at most \(1\) for \(k\ge7\).
\end{problem}

\begin{problem}\label{pr:sublattices}
For every integer \(n'\ge1\), \(\Z^2\) has \(\sum_{e\mid n'}e\) sublattices of index \(n'\) (\(3\) and \(7\) for \(n'=2,4\)); hence \(\sum_{L'}[\Z^2:L']^{-\omega}=\zeta(\omega)\zeta(\omega-1)\) for real \(\omega>2\), the sum over all sublattices \(L'\) of finite index.
\end{problem}

Recall from Chapter~\ref{m1-chap} that \(\Dmain\le(M/b^6)^{kf}\Did\) for the data; and \(kf\ge40\), since \(k\ge1\) and \(f\ge x^{11}\varphi(q_0)\ge40\) (Problems~\ref{da-pr:m-big} and~\ref{da-pr:T}).

\begin{problem}\label{pr:exponent}
With \(z_0:=p_j^{2f}\), \((b^{2f}/z_0)^k\Dmain\le b^{-3kf}\Did\le b^{-30}\Did\).
\end{problem}

\begin{problem}\label{pr:strip-sum}
The second moment of \(U(Y_0)(F)\smallsetminus R(F)\) over \(\Bc_j\) has square root at most \(2b^{-15}\sqrt{\Did}\), and \(4b^{-30}\le b^{-12}\).
\end{problem}

Problems~\ref{pr:scalar-one}, \ref{pr:strip-meets}, \ref{pr:hs-blocks}, \ref{pr:exponent} and~\ref{pr:strip-sum}, with Theorems~\ref{thm:link-bound} and~\ref{thm:cut}, prove the following.

\begin{theorem}[the strip]\label{thm:strip}
For the data and \(j=0,1\),
\[
 \sum_{s\in\Bc_j}\lVert\Kc_s-\rB(s)\rVert^2\le b^{-12}\Did .
\]
\end{theorem}

\noindent The proof states Theorem~\ref{thm:strip} for an admissible filling over a CM field equipped with a CM type and containing a primitive \(5040\)-th root of unity, any representation \(\rho\) of its finite group on a finite-dimensional complex inner product space, unitary or not, and any linear map from the coordinate space that satisfies the block identity of Theorem~\ref{m1-thm:blocks}.  The data satisfy these hypotheses.

Recall Problem~\ref{m1-pr:near-one}: under the hypotheses of Problem~\ref{m1-pr:bessel-translate}, which hold for the data, if \(\sum_{t\in\Rb}\lVert\Kc_t-\rB(t)\rVert^2\le b^{-12}\Did\), then
\[
 \lVert\Pc^{\circ*}\Pc^\circ-1\rVert^2\le b^{-12}D_X ;
\]
and Problem~\ref{m1-pr:P-one}: then \(\lVert\Kc_1\rVert^2\le2D_X\) for \(b\ge1000\).

\begin{problem}\label{pr:near-one-data}
For the data, \(\lVert\Pc^{\circ*}\Pc^\circ-1\rVert^2\le b^{-12}D_X\) and \(\lVert\Kc_1\rVert^2\le2D_X\).
\end{problem}

\paragraph{What has been proved.}  Theorem~\ref{thm:strip} is the bound \(b^{-12}\Did\) for the sum of the squared Hilbert--Schmidt errors on each strip, and the first inequality of Problem~\ref{pr:near-one-data} makes \(\Pc^\circ\) a near isometry in the sense of Chapter~\ref{bp-chap}.  Taken as given: the block identity and the Bessel functions (Theorems~\ref{m1-thm:blocks} and~\ref{m1-thm:bessel-system}), and Theorems~\ref{thm:survivors}, \ref{thm:link-bound} and~\ref{thm:cut} and the inequality \(\Dmain\le(M/b^6)^{kf}\Did\).  Problem~\ref{pr:strip-meets} gives a lower bound for the meet norm on each strip only.  Chapters~\ref{m3-chap} and~\ref{m4-chap} use the open reductions of Definition~\ref{def:open} and the fixers of Definition~\ref{def:fixer} to bound \(\sum_{s\in\Gb}\lVert\Kc_s\rVert^2\) by \((1+9b^{-4})\Dmain\) (Theorem~\ref{m4-thm:second-moment}).

\begin{leanremark}
\raggedright In the Lean proof \cite{GL}: Theorem~\ref{thm:strip} is \tl{Filling.HasAmplitudeBlocks.}\allowbreak\tl{sum\_hilbertSchmidtNorm\_sub\_besselTranslate\_sq\_le}; Theorem~\ref{thm:survivors} is \tl{Filling.AvoidsScalarCosets.amplitudeSum\_compl\_besselS\_eq}; Theorem~\ref{thm:link-bound} is \tl{Filling.linkSecondMoment\_le\_of\_le\_meetNorm}; Theorem~\ref{thm:cut} is \tl{Filling.sqrt\_secondMoment\_compl\_besselS\_le}.  The discriminant, the link, the meet norm, open reductions, fixers and the avoidance of scalar cosets are \tl{Filling.termDisc}, \tl{link}, \tl{Filling.meetNorm}, \tl{Filling.IsOpenElement}, \tl{Filling.fixer} and \tl{Filling.AvoidsScalarCosets}.  The pieces of a cut block are \tl{Filling.piece}.  Hua's integral, the sublattice sum and Liouville's inequality in the proof of Theorem~\ref{thm:link-bound} are \tl{lintegral\_norm\_det\_add\_inv\_pow\_le}, \tl{Filling.sum\_discriminant\_inv\_le} and \tl{Filling.liouville\_gram}.
\end{leanremark}

\section*{Hints}

\subsection*{\S\ref{sec:characters}}
\begin{description}
\item[\ref{pr:agree}] Evaluate \(\phi(ugu')\) in two ways.
\item[\ref{pr:gl2-cells}] For \(g=\diag(t_1,t_2)\) and \(u'=\left(\begin{smallmatrix}1&y\\0&1\end{smallmatrix}\right)\), find the \(u\) with \(ugu'=g\).  In the big cell \(U\cap wUw^{-1}=1\).
\item[\ref{pr:gl3-cells}] Sage or Python: run over \(\GL_3(\F_p)\), form each \(UgU\) with the value \(\psi_U(u)\psi_U(u')\) attached to \(ugu'\), and test consistency.
\item[\ref{pr:gl2-frame}] Use Problem~\ref{pr:gl2-cells}; the \(g\) with \(g_{21}=0\) form the upper triangular group.
\end{description}

\subsection*{\S\ref{sec:frames}}
\begin{description}
\item[\ref{pr:two-links}] Compute the matrices of frame products of \(\gamma_\pm\), and the intersection \(\Uc\cap\gamma_\pm\Uc\gamma_\pm^{-1}\).
\item[\ref{pr:scalar-one}] If \(\zeta\ne1\), the ideal \((\zeta-1)\) is divisible by the \(f\) distinct primes \(w\in\cMc_j\), of norm \(p_j\); bound \(\lvert\sigma(\zeta)-1\rvert\le2\) for every embedding \(\sigma\).
\item[\ref{pr:strip-meets}] An integral element \(t\) of \(R(\A_f)\) fixes \(v_1\) and satisfies \(h(ty,v_1)=h(y,v_1)\); reduce modulo \(w\), and deduce that the adjoint inverse of \(t_w\) also fixes \(v_1\).  For a scalar coset, compare the entries at \(v_1\) of \(s_w\) and \(t_w\) and use Problem~\ref{pr:scalar-one}.  For \(p_j^2\ge M/b\) use \(p_{1-j}\le b\,p_j\).
\item[\ref{pr:fixer-meets}] Which vectors does \(U(X_0)\) fix?  Compare the definition of a fixer with Definition~\ref{def:meet}.
\end{description}

\subsection*{\S\ref{sec:gain}}
\begin{description}
\item[\ref{pr:petersson-gain}] \(p\mid\gamma_{21}\); then Problem~\ref{ps-pr:J-bound}.
\item[\ref{pr:hs-blocks}] Expand the norm in the basis \(e_{c,i}\), and split \(U(Y_0)(F)\) into \(R(F)\) and its complement.
\item[\ref{pr:cross}] Expand the square.
\item[\ref{pr:hua}] Polar coordinates: the sphere \(S^{2k-1}\) has volume \(2\pi^k/(k-1)!\).
\item[\ref{pr:sublattices}] Each sublattice is spanned by the columns of exactly one \(\left(\begin{smallmatrix}g_1&g_3\\0&g_2\end{smallmatrix}\right)\) with \(g_1g_2=n'\), \(0\le g_3<g_1\).
\item[\ref{pr:exponent}] Problem~\ref{pr:strip-meets}.
\item[\ref{pr:strip-sum}] Theorem~\ref{thm:cut} with \(\Bc=\Bc_j\) and \(C_0=b^{-30}\Did\), after Theorem~\ref{thm:link-bound} and Problems~\ref{pr:strip-meets} and~\ref{pr:exponent}.
\item[\ref{pr:near-one-data}] \(\Rb\subseteq\Bc_0\).
\end{description}

\begin{thebibliography}{GL}
\bibitem[Ca72]{Ca72} R.~W.~Carter, \emph{Simple Groups of Lie Type}, Wiley, London, 1972.
\bibitem[Ca85]{Ca85} R.~W.~Carter, \emph{Finite Groups of Lie Type: Conjugacy Classes and Complex Characters}, Wiley, Chichester, 1985.
\bibitem[GL]{GL} L.~Sernau, \emph{GRH for Dirichlet \(L\)-functions, in Lean~4}, a shorter rewriting of the Lean proof, text by Claude, \url{https://github.com/lsernau/grh-lean}, commit \texttt{72fead1}, readable edition \texttt{grh/}.
\bibitem[Hua]{Hua} L.-K.~Hua, \emph{Harmonic Analysis of Functions of Several Complex Variables in the Classical Domains}, Transl. Math. Monographs \textbf{6}, Amer. Math. Soc., Providence, 1963.
\bibitem[Sch]{Sch} W.~M.~Schmidt, \emph{Asymptotic formulae for point lattices of bounded determinant and subspaces of bounded height}, Duke Math. J. \textbf{35} (1968), 327--339.
\end{thebibliography}

\end{document}
